\documentclass[12pt]{article}
\usepackage[utf8]{inputenc}
\usepackage[T1]{fontenc}
\usepackage{hyperref}
\title{Lawunuia Language Resources for K-12 Curriculum Development}
\author{Compiled by Claude Sonnet 5 (claude-sonnet-5) for Andr\'as Kornai}
\date{}
\begin{document}
\maketitle

\section{Introduction}
Lawunuia (ISO 639-3 tgi; also called Nagarege, Nagarige, Naghareghe and Piva) is an Austronesian language. The search results place it along the Piva River in the Autonomous Region of Bougainville, Papua New Guinea. The results do not support any link to Milne Bay Province. It is closely related to Banoni. Very little documentation exists.

\section{Primary Educational Resources}
No primers, textbooks or school programmes were found.

\section{Linguistic and Reference Materials}
\begin{itemize}
\item Wikipedia overview: \url{https://en.wikipedia.org/wiki/Lawunuia_language}
\item Counting-system data for Lawunuia: \url{https://lingweb.eva.mpg.de/channumerals/Lawunuia.htm}
\item History of research on the Austronesian languages of Milne Bay Province (ANU). This surfaced in the search but concerns Milne Bay and probably does not cover Lawunuia: \url{https://openresearch-repository.anu.edu.au/items/9966de95-7d52-4a35-a30a-3b550f59e8e1}
\end{itemize}

\section{Religious and Community Resources}
None found.

\section{Digital and Audio Resources}
\begin{itemize}
\item Endangered Languages Project record: \url{https://endangeredlanguages.com/node/86628}
\item Verbix language page: \url{https://docs.verbix.com/Languages/Lawunuia}
\end{itemize}
A search snippet mentions that PARADISEC holds 1973 Banoni/Piva recordings. No URL was seen for them.

\section{Curriculum Development Context}
There is almost nothing published for Lawunuia. A curriculum would need fieldwork and community partnership, possibly building on documentation of the related Banoni language. Estimates of speaker numbers differ across sources.

\end{document}